\documentclass[11pt,a4paper]{letter}
\usepackage[T1]{fontenc}
\usepackage{lmodern}
\usepackage[utf8]{inputenc}
\usepackage[margin=22mm]{geometry}
\usepackage{enumitem}
\usepackage[hidelinks]{hyperref}
\pagestyle{empty}
\date{13 September 2026}

\signature{Yeoneung Kim\\
Department of Industrial Engineering, Yonsei University\\
\href{mailto:yeoneung@yonsei.ac.kr}{yeoneung@yonsei.ac.kr}}
\address{Department of Industrial Engineering, Yonsei University\\
50 Yonsei-ro, Seodaemun-gu\\
Seoul 03722, Republic of Korea}

\begin{document}
\begin{letter}{Editors\\Simulation Modelling Practice and Theory}
\opening{Dear Editors,}

Please consider the manuscript ``Simulation-based comparison of offline value
learning and online mixed-integer control for battery dispatch'' as a research
article in \emph{Simulation Modelling Practice and Theory}.

The paper uses a common simulation model to compare offline value learning
with online mixed-integer control. A discontinuous import-band fee makes
accurate deployment of the billed objective essential. The comparison
separates the contribution of learning from the effects of action search,
information, and computational budgets. The main contributions are:

\begin{itemize}[leftmargin=1.5em,itemsep=5pt,topsep=4pt,parsep=0pt]
\item \textbf{A matched simulation benchmark.} Offline learning and
conditional-mean and two-stage stochastic MIQPs share the economic objective,
information, terminal cost, feasible actions, and 15-minute update interval.
Nested scenario sets and explicit solve budgets distinguish stochastic-model
resolution from the performance of the implemented optimizer.

\item \textbf{Deployment-consistent learning and attribution.} Smooth
finite-horizon value learning is coupled to hard-cost one-step minimization
over the deployed feasible set. A matched reference-only controller retains
the same decision rule while removing the neural correction, separating the
economic contribution of learning from that of the one-step search.

\item \textbf{Numerical validation and computational trade-offs.} Finite-horizon
Bellman bounds separate residual variation, action-search, expectation, and
terminal errors. Four seasons, repeated training batches, and tariff
sensitivity support a joint assessment of operating cost, online latency,
and the deployment period needed to amortize offline computation.
\end{itemize}

The learned correction reduces cost relative to its no-learning counterpart
by 2.9\%, 11.9\%, and 24.6\% at the three fee anchors, while the primary
two-scenario MIQP achieves lower operating cost. Together with measured
online and offline computation, these results quantify the trade-off between
operating cost and controller latency. The focus on simulation-model
implementation, numerical validation, and controlled comparison of decision
procedures fits the journal's readership.

Source data are public, and the code, stored weights, and numerical
outputs are available in the GitHub repository at
\url{https://github.com/yeoneung/rvpinn}; the manuscript identifies the exact
reproducibility version.

\closing{Sincerely,}
\end{letter}
\end{document}
